% !TEX root = ../main.tex
\chapter{おわりに}\label{chap:conclusion}
% ===== 11/26 書き直し(Claude)。レオが読んで理解 → 確認の問い → 直す =====

\section{まとめ}
本研究では、Raspberry Pi 4B のマイクロクラスタを対象に、空間熱干渉を外部熱影響 $H$ として状態に持つ熱モデルとカルマンフィルタで 30 秒先の温度と不確かさを予測し、予測温度が安全閾値を超えるときだけファンを回す制御を扱った。対外発表の評価の問題を直し、3 台と 8 台の実機で評価した。得られた結果は次のとおりである。

\begin{enumerate}
  \item 空間熱干渉:1 台の排熱が隣を直接温める干渉は、囲いの有無・位置関係によらず 1\,℃ 未満であった。熱は囲いの中に溜まる形で効き、囲うだけで 5〜10\,℃ 高くなった。隣のファンの風は周りを 2〜3\,℃ 冷やした
  \item 予測:$H$ を持つモデルは、30 秒先の予測誤差を約半分にした(囲いなし 0.64 → 0.39\,℃、囲いあり 0.84 → 0.44\,℃)。囲いありでは、囲いの中に温度計を置いたのと同じ精度を、温度計なしで達成した
  \item ファン制御:囲いのないラックで、80\,℃ を超えずに稼働を常時 ON の約 4 割にした。ただし、しきい値方式との差はほぼなかった。囲いの中ではファンの余力が小さく\TBD{11 月 26 日の結果}
\end{enumerate}

これらの結果から、空間熱干渉は隣の 1 台からの排熱ではなく、囲いに溜まる熱として現れ、$H$ はそれを温度計なしで推定できる。一方、速く強いファンでは予測としきい値の差は小さく、予測が要るのは、効きの遅い・弱い冷却を用いる場合であると考えられる。

本研究では、空間熱干渉を外部熱影響として推定することで、干渉のある環境でも追加の温度計なしに温度を精度よく予測できること、およびその予測に基づくファン制御により、80\,℃ を超えずにファンの稼働時間を常時 ON の約 4 割に減らせること(囲いのない環境)を実機で示した。
% 11/26 レオ:題目どおり、達成は 2 つ(干渉のある環境で予測ができる/それでファンの稼働を減らせる)。囲いありの制御の結果が良ければ、括弧の但し書きを外す

\section{今後の課題}
本論文の結果は予備審査の時点のものである。11 月 26 日から次の実験を進めており、12 月の追加実験と合わせて最終版(1 月)に反映する。

\begin{itemize}
  \item ファンの段階を予測で選ぶ制御:段階ごとに 30 秒先の温度を予測し、80\,℃ を守れる一番弱い段階を選ぶ。しきい値や OS 標準と、稼働と静かさで比べる
  \item 負荷のかけ方:一気にかける形と徐々にかける形、熱い場所が移っていく形で比べる。熱い場所が動く形は、自ノードの負荷は同じで周りの熱だけが変わるため、$H$ の効果を制御で確かめる条件になる
  \item 比べる相手:1 分おき・ランダムにファンを回す単純な方式、OS 標準のファン制御
  \item 効きの遅い冷却:部屋のサーキュレーターなど、効きの遅い冷却手段で、シミュレーションで得た境界(むだ時間 15〜30 秒)を実機で確かめる
  \item 囲いの大きさ・開口、実際のアプリケーションの負荷、ファンの風の向きの確認
\end{itemize}
